\chapter{การทบทวนวรรณกรรมที่เกี่ยวข้อง}
\label{chapter:literature-review}

ในการพัฒนา UniAssist AI ระบบผู้ช่วยอัจฉริยะเพื่อการเรียนรู้และการติดตามนักศึกษา คณะผู้จัดทำได้ศึกษาทฤษฎีและบทความวิจัยที่เกี่ยวข้อง เพื่อให้ระบบสามารถทำงานได้อย่างมีประสิทธิภาพ แม่นยำ และตอบโจทย์ผู้ใช้งาน โดยแบ่งหัวข้อการศึกษาออกเป็นสองส่วน ได้แก่ ทฤษฎีที่เกี่ยวข้อง และบทความวิจัยที่เกี่ยวข้อง

\section{ทฤษฎีที่เกี่ยวข้อง}

\subsection{โมเดลภาษาขนาดใหญ่ (Large Language Model: LLM)}

โมเดลภาษาขนาดใหญ่คือโมเดลปัญญาประดิษฐ์ประเภทการเรียนรู้เชิงลึก (Deep Learning) ที่ได้รับการฝึกฝนด้วยข้อมูลข้อความจำนวนมหาศาล โดยมีพื้นฐานมาจากสถาปัตยกรรมแบบ Transformer ซึ่งอาศัยกลไก Self-Attention ในการเรียนรู้ความสัมพันธ์ระหว่างคำในบริบท \cite{vaswani2017attention} ทำให้มีความสามารถในการเข้าใจบริบททางภาษา การสรุปความ และการสร้างข้อความที่ใกล้เคียงกับมนุษย์ ในโครงงานนี้มีการศึกษาเรื่องการออกแบบคำสั่ง (Prompt Engineering) และการให้เหตุผลของโมเดล (LLM Reasoning) เพื่อนำมาประยุกต์ใช้ในการวิเคราะห์เจตนาของคำถาม (Intent Recognition) การแปลงคำถามเป็นคำสั่ง SQL และการให้เหตุผลเกี่ยวกับกฎระเบียบการลงทะเบียนเรียนที่ซับซ้อน

\subsection{การแปลงภาษาธรรมชาติเป็นคำสั่ง SQL (Text-to-SQL)}

Text-to-SQL หรือ NL2SQL คือกระบวนการแปลงคำถามภาษาธรรมชาติให้เป็นคำสั่งภาษา SQL โดยอัตโนมัติ เพื่อให้ผู้ใช้ที่ไม่มีความรู้ด้านฐานข้อมูลสามารถสอบถามข้อมูลจากฐานข้อมูลเชิงสัมพันธ์ได้ กระบวนการนี้ต้องอาศัยความเข้าใจทั้งความหมายของคำถาม (Semantics) และโครงสร้างของฐานข้อมูล (Schema) ในอดีตงานด้านนี้ใช้แบบจำลองที่ฝึกเฉพาะทางบนชุดข้อมูลมาตรฐาน เช่น WikiSQL และ Spider \cite{yu2018spider} แต่ในปัจจุบันโมเดลภาษาขนาดใหญ่ที่มีความสามารถในการให้เหตุผลสามารถสร้างคำสั่ง SQL ได้อย่างมีประสิทธิภาพผ่านการให้บริบทของ schema และตัวอย่างในคำสั่ง (In-context Learning) โดยไม่จำเป็นต้องฝึกโมเดลใหม่ ข้อได้เปรียบของแนวทาง Text-to-SQL เมื่อเทียบกับการให้ LLM ตอบจากความจำโดยตรง คือคำตอบจะถูกดึงจากข้อมูลจริงในฐานข้อมูล ทำให้ตรวจสอบย้อนกลับได้และลดปัญหาการสร้างข้อมูลเท็จ เนื่องจากข้อมูลหลักสูตรมีลักษณะเป็นข้อมูลที่มีโครงสร้างชัดเจน แนวทางนี้จึงมีความเหมาะสมสำหรับการค้นคืนข้อมูลเชิงข้อเท็จจริง เช่น รายวิชา หน่วยกิต และเงื่อนไขบังคับก่อน

\subsection{การสร้างข้อความโดยอาศัยการค้นคืนข้อมูล (Retrieval-Augmented Generation: RAG)}

RAG \cite{lewis2020rag} คือกระบวนการที่ช่วยเพิ่มความแม่นยำให้กับโมเดลภาษา โดยการเชื่อมต่อโมเดลเข้ากับฐานความรู้ภายนอก (External Knowledge Base) เช่น คู่มือนักศึกษา หรือระเบียบการลงทะเบียน กระบวนการทำงานประกอบด้วย 3 ส่วนหลัก คือ (1) Retrieve ค้นหาเอกสารที่เกี่ยวข้องกับคำถามจากฐานข้อมูล (2) Augment นำข้อมูลที่ค้นพบมารวมกับคำถามตั้งต้น และ (3) Generate ส่งข้อมูลทั้งหมดให้ LLM ประมวลผลเพื่อสร้างคำตอบ เทคนิคนี้ช่วยลดปัญหาการตอบข้อมูลเท็จ (Hallucination) และทำให้ข้อมูลมีความเป็นปัจจุบัน ในโครงงานนี้แนวคิด RAG ถูกนำมาใช้เป็นกรอบความคิดเชิงหลักการของการ ``ตอบคำถามโดยยึดจากข้อมูลจริง'' โดยแนวทาง Text-to-SQL ที่นำมาใช้ถือเป็นรูปแบบหนึ่งของการค้นคืนข้อมูลที่เหมาะกับข้อมูลเชิงโครงสร้าง ขณะที่โหมดสนทนาทั่วไปของระบบใช้การยึดข้อมูล (Grounding) จากสรุปข้อเท็จจริงในฐานข้อมูลเพื่อลดการตอบนอกบริบท

\subsection{ฐานข้อมูลเชิงเวกเตอร์และการฝังคำ (Vector Database \& Text Embedding)}

ทฤษฎีการแปลงข้อมูลข้อความให้อยู่ในรูปแบบของตัวเลขหรือเวกเตอร์ (Vector Embeddings) เพื่อให้คอมพิวเตอร์สามารถเข้าใจความหมายเชิงบริบท (Semantic Meaning) ของข้อความได้ การใช้งานฐานข้อมูลเชิงเวกเตอร์ช่วยให้ระบบสามารถค้นหาข้อมูลที่มีความหมายใกล้เคียงกัน (Semantic Search) ได้อย่างมีประสิทธิภาพ แม้ว่าผู้ใช้งานจะใช้คำศัพท์ที่ไม่ตรงกับเอกสารต้นฉบับก็ตาม แนวคิดนี้เหมาะกับข้อมูลที่ไม่มีโครงสร้างชัดเจน (เช่น ประกาศหรือคู่มือแบบข้อความยาว) และเป็นแนวทางที่โครงงานพิจารณาไว้สำหรับการต่อยอดในการค้นคืนข้อมูลกฎระเบียบเชิงบรรยายในอนาคต

\subsection{ระบบแนะนำข้อมูล (Recommender System)}

การศึกษาระบบแนะนำเพื่อนำมาประยุกต์ใช้ในการแนะนำรายวิชา (Course Recommendation) โดยในโครงงานนี้มุ่งเน้นศึกษาระบบแนะนำที่อิงตามความรู้และข้อจำกัด (Knowledge-based / Constraint-based Recommender System) เนื่องจากเงื่อนไขการลงทะเบียนเรียนมีความซับซ้อนและมีกฎเกณฑ์ตายตัว (เช่น รายวิชาบังคับก่อน หรือ Prerequisite) ระบบจึงต้องวิเคราะห์ความต้องการของผู้เรียนร่วมกับโครงสร้างหลักสูตรเพื่อเสนอทางเลือกที่เหมาะสมที่สุด

\subsection{การทำเหมืองข้อมูลทางการศึกษา (Educational Data Mining: EDM)}

EDM คือศาสตร์ที่ประยุกต์ใช้เทคนิคการทำเหมืองข้อมูล (Data Mining) กับข้อมูลทางการศึกษา เพื่อค้นหารูปแบบความสัมพันธ์ที่ซ่อนอยู่ ในโครงงานนี้ได้ศึกษาเทคนิคการวิเคราะห์ข้อมูลผลการเรียน (Grade Analysis) และการประเมินความเสี่ยง (Risk Assessment) โดยใช้การกำหนดกฎเกณฑ์ (Rule-based) จากระเบียบข้อบังคับของสถาบัน เพื่อจำแนกและแจ้งเตือนแนวโน้มที่นักศึกษาจะประสบปัญหาทางการเรียน เช่น การติดสถานะภาคทัณฑ์ (Probation) หรือการพ้นสภาพ (Dismissal)

\subsection{สถาปัตยกรรมซอฟต์แวร์และการพัฒนาเว็บ (Software Architecture \& Web Development)}

การศึกษาโครงสร้างการพัฒนาซอฟต์แวร์ในรูปแบบเว็บแอปพลิเคชันด้วยเฟรมเวิร์ก Flask และการออกแบบส่วนติดต่อกับบริการ LLM ผ่านมาตรฐาน API ที่เข้ากันได้กับ OpenAI (OpenAI-compatible API) รวมถึงการออกแบบฐานข้อมูลเชิงสัมพันธ์ (Relational Database) ด้วย SQLite เพื่อจัดเก็บข้อมูลหลักสูตรและผลการเรียนให้มีความถูกต้อง (Data Integrity) และปลอดภัย ตลอดจนการยืนยันตัวตนผู้ใช้ด้วยมาตรฐาน OAuth 2.0 เพื่อควบคุมสิทธิ์การเข้าถึงตามบทบาทของผู้ใช้

\section{บทความวิจัยที่เกี่ยวข้อง}

\subsection{Systematic Analysis of Retrieval-Augmented Generation-Based LLMs for Medical Chatbot Applications \cite{medicalRAG}}

งานวิจัยนี้มีวัตถุประสงค์เพื่อแก้ปัญหาการให้ข้อมูลที่ผิดพลาด (Hallucinations) ของโมเดลภาษาขนาดใหญ่ในบริบททางการแพทย์ ซึ่งความถูกต้องแม่นยำและความโปร่งใสของข้อมูลเป็นสิ่งสำคัญสูงสุด โดยมุ่งเน้นการพัฒนาระบบที่สามารถทำงานได้ในสภาพแวดล้อมที่มีทรัพยากรการคำนวณจำกัด ผู้วิจัยได้นำเสนอการเปรียบเทียบสถาปัตยกรรม RAG 3 รูปแบบ ได้แก่ Naive RAG, Advanced RAG และ Modular RAG พบว่ารูปแบบ Advanced RAG มีความเหมาะสมที่สุดสำหรับการพัฒนาแชตบอตทางการแพทย์ เนื่องจากมีความสามารถในการจัดการบริบท (Context-awareness) ที่ซับซ้อนได้ดีกว่า ผลการทดลองยังพบว่าโมเดล Mistral-7B ที่ผ่านการปรับจูนและใช้งานร่วมกับ RAG ให้ประสิทธิภาพสูงสุด และการใช้ RAG ช่วยปรับปรุงความถูกต้องของข้อมูลได้อย่างมีนัยสำคัญเมื่อเทียบกับโมเดลที่ปรับจูนเพียงอย่างเดียว งานวิจัยนี้ชี้ให้เห็นความสำคัญของการยึดคำตอบกับแหล่งข้อมูลจริง ซึ่งเป็นหลักการเดียวกับที่โครงงานนี้นำมาใช้ผ่านแนวทาง Text-to-SQL

\subsection{Development of an Academic Services Chatbot Based on Retrieval-Augmented Generation (RAG) \cite{academicChatbotRAG}}

งานวิจัยนำเสนอการพัฒนาระบบแชตบอทสำหรับการให้บริการข้อมูลทางวิชาการในระดับมหาวิทยาลัย เพื่อแก้ปัญหาการให้บริการข้อมูลแบบดั้งเดิมที่ต้องอาศัยเจ้าหน้าที่เป็นหลัก ซึ่งก่อให้เกิดความล่าช้าและภาระงานซ้ำซ้อน ระบบใช้โมเดลภาษาขนาดใหญ่เป็นแกนหลักร่วมกับแนวคิด RAG เพื่อแก้ข้อจำกัดเรื่อง hallucination และการเข้าถึงข้อมูลเฉพาะโดเมน โดยรวบรวมเอกสารและข้อมูลทางวิชาการภายในสถาบัน แบ่งออกเป็นหน่วยย่อยและจัดเก็บในฐานข้อมูลเวกเตอร์ พร้อมทั้งนำแนวคิด Conversational Memory มาใช้เพื่อรองรับการสนทนาแบบต่อเนื่อง ผลการประเมินพบว่าระบบสามารถตอบคำถามเชิงข้อเท็จจริงได้อย่างมีประสิทธิภาพ แต่ยังมีข้อจำกัดในการตอบคำถามเชิงการให้เหตุผลที่ซับซ้อน งานวิจัยนี้เป็นตัวอย่างของการนำแชตบอทมาใช้ในบริบทการให้บริการวิชาการซึ่งตรงกับกลุ่มเป้าหมายของโครงงาน

\subsection{Towards Effective Teaching Assistants: From Intent-based Chatbots to LLM-powered Teaching Assistants \cite{teachingAssistantsLLM}}

งานวิจัยนี้เปรียบเทียบประสิทธิภาพระหว่างแชตบอตแบบดั้งเดิมที่ใช้การกำหนดเจตนา (Intent-based) กับแชตบอตยุคใหม่ที่ขับเคลื่อนด้วยโมเดลภาษาขนาดใหญ่ร่วมกับเทคนิค RAG เพื่อใช้เป็นผู้ช่วยสอน (Teaching Assistant) ในระยะแรกผู้วิจัยพัฒนาแชตบอตแบบ Intent-based ด้วยแพลตฟอร์ม Amazon Alexa Skills Kit ซึ่งพบข้อจำกัดคือการสร้างฐานความรู้ด้วยมือมีความล่าช้าและขาดความยืดหยุ่นในการตอบคำถามนอกรูปแบบที่กำหนด ในระยะที่สองจึงเปลี่ยนมาใช้แนวทาง RAG ด้วยเฟรมเวิร์ก LangChain ร่วมกับโมเดล GPT-3.5 Turbo ผลการเปรียบเทียบพบว่าแนวทาง RAG เหนือกว่าอย่างชัดเจนทั้งด้านความยืดหยุ่นและความสามารถในการปรับตัว โดยได้คะแนนความถูกต้องและความมีประโยชน์โดยรวมในระดับสูง งานวิจัยนี้ยืนยันข้อได้เปรียบของการใช้ LLM แทนระบบที่อิงรูปแบบประโยคตายตัว ซึ่งเป็นเหตุผลที่โครงงานเลือกใช้ LLM ในการตีความคำถามที่หลากหลาย

\subsection{RAG-Based AI Chatbot for Student and Institutional Assistance \cite{ragStudentInstitutional}}

งานวิจัยนำเสนอการพัฒนาระบบแชตบอทเพื่อสนับสนุนการให้บริการข้อมูลทางวิชาการและการบริการภายในสถาบันอุดมศึกษา เพื่อแก้ปัญหาการให้บริการข้อมูลแบบดั้งเดิมที่ต้องพึ่งพาเจ้าหน้าที่ โดยชี้ให้เห็นว่าคำถามจากนักศึกษาที่เกี่ยวกับการรับสมัคร การเรียนการสอน และระเบียบการศึกษามีจำนวนมากและมีลักษณะซ้ำกัน ระบบใช้ LLM เป็นแกนหลักร่วมกับ RAG แบบหลายขั้นตอน (multi-stage pipeline) เพื่อเพิ่มคุณภาพของบริบท พร้อมทั้งนำ Conversational Context มาใช้รองรับการอ้างอิงคำถามก่อนหน้า ผลการทดลองแสดงให้เห็นว่าระบบสามารถตอบคำถามเชิงข้อเท็จจริงและจัดการคำถามนอกขอบเขตความรู้ได้อย่างเหมาะสม แต่ยังพบข้อจำกัดในการให้เหตุผลที่ซับซ้อน ซึ่งเป็นความท้าทายที่โครงงานนี้ให้ความสำคัญเช่นกัน

\subsection{OpenThaiGPT 1.5: A Thai-Centric Open Source Large Language Model \cite{openthaigpt2024}}

รายงานทางเทคนิคนี้เป็นการพัฒนาโมเดลภาษาขนาดใหญ่ที่มีความเชี่ยวชาญเฉพาะสำหรับภาษาไทย โดยพัฒนาต่อยอดจากสถาปัตยกรรมของโมเดล Qwen v2.5 มีให้เลือกใช้ 2 ขนาด คือ 7B และ 72B พารามิเตอร์ ในกระบวนการพัฒนาใช้การปรับจูนตามคำสั่ง (Instruction Finetuning) ด้วยเทคนิค LoRA บนชุดข้อมูลคุณภาพสูงกว่า 2 ล้านคู่คำสั่ง จุดเด่นสำคัญคือการรองรับการสนทนาแบบต่อเนื่อง การทำงานร่วมกับ RAG และความสามารถในการเรียกใช้เครื่องมือ (Tool Calling) ผลการทดสอบบนชุดข้อสอบระดับชาติของไทยพบว่ารุ่น 72B ทำคะแนนเฉลี่ยรวมได้สูงกว่าโมเดลโอเพนซอร์สอื่นในระดับเดียวกัน งานวิจัยนี้เป็นข้อมูลประกอบการพิจารณาเลือกโมเดลภาษาที่รองรับภาษาไทยได้ดีสำหรับระบบที่ผู้ใช้สื่อสารเป็นภาษาไทย

\subsection{Qwen2.5 Technical Report \cite{qwen2025technical}}

รายงานทางเทคนิค Qwen2.5 นำเสนอการพัฒนาโมเดลภาษาขนาดใหญ่รุ่นล่าสุดของตระกูล Qwen ที่มุ่งยกระดับประสิทธิภาพด้านความเข้าใจภาษา การให้เหตุผล คณิตศาสตร์ การเขียนโค้ด และการปฏิบัติตามคำสั่งของผู้ใช้ ด้านสถาปัตยกรรมยังคงใช้โครงสร้าง Transformer-based decoder ที่ปรับปรุงด้วยเทคนิค Grouped Query Attention (GQA), SwiGLU, Rotary Positional Embedding (RoPE) และ RMSNorm โมเดลมีให้เลือกหลายขนาดตั้งแต่ 0.5B ถึง 72B พารามิเตอร์ กระบวนการฝึกเน้นการขยายคุณภาพและปริมาณข้อมูลถึง 18 ล้านล้านโทเคน จุดเด่นสำคัญคือการรองรับการสร้างข้อความยาว การจัดการข้อมูลเชิงโครงสร้าง (เช่น ตารางและ JSON) และความสามารถในการเรียกใช้เครื่องมือ ซึ่งเป็นคุณสมบัติที่เหมาะสมอย่างยิ่งต่อการสกัดข้อมูลหลักสูตรให้อยู่ในรูปแบบ JSON และการสร้างคำสั่ง SQL ในโครงงานนี้ จึงเป็นเหตุผลหลักที่โครงงานเลือกใช้โมเดลตระกูล Qwen เป็นโมเดลหลักของระบบ
